% ============================================================================
% Font setup for the open-source edition
% Compatible with ElegantBook (ctex already loaded by elegantbook.cls)
% ============================================================================

% ============================================================================
% Local font directory
% ============================================================================
\newcommand{\fontdir}{assets/fonts/}

% ============================================================================
% CJK font setup (xeCJK already loaded by ctex)
% ============================================================================
\xeCJKsetup{
    AutoFakeBold = true,
    AutoFakeSlant = 0.2,
    PunctStyle = kaiming,
    CheckSingle = true,
    CJKglue = \hskip 0pt plus 0.04em minus 0.02em,
    CJKecglue = \hskip 0.25em plus 0.08em minus 0.05em
}

% ============================================================================
% Latin fonts
% ============================================================================
\setmainfont{TimesNewRoman.ttf}[
    Path = \fontdir,
    Scale = 1.0,
    Ligatures = TeX,
    BoldFont = TimesNewRomanBold.ttf,
    ItalicFont = TimesNewRomanItalic.ttf,
    BoldItalicFont = TimesNewRomanBoldItalic.ttf,
]

\setsansfont{GoogleSans-Regular.ttf}[
    Path = \fontdir,
    Scale = 0.95,
    BoldFont = GoogleSans-Bold.ttf,
    ItalicFont = GoogleSans-Italic.ttf,
    BoldItalicFont = GoogleSans-BoldItalic.ttf,
]

\setmonofont{CourierNew.ttf}[
    Path = \fontdir,
    Scale = 0.9,
    BoldFont = CourierNewBold.ttf,
    ItalicFont = CourierNewItalic.ttf,
    BoldItalicFont = CourierNewBoldItalic.ttf,
]

% ============================================================================
% CJK font commands
% ============================================================================
% 思源宋体（NotoSerifSC，Regular/Bold 双字重）作正文宋体；
% 思源黑体（NotoSansSC，Regular/Bold 双字重）作黑体，与思源宋体同族配套。
% 楷体/仿宋沿用 STKaiti/STFangsong，单字重，粗体由 AutoFakeBold 合成。
\newcommand{\MySongTi}{NotoSerifSC-Regular.otf}
\newcommand{\MySongTiBold}{NotoSerifSC-Bold.otf}
\newcommand{\MyHeiTi}{NotoSansSC-Regular.otf}
\newcommand{\MyHeiTiBold}{NotoSansSC-Bold.otf}
\newcommand{\MyKaiTi}{STKaiti.ttf}
\newcommand{\MyFangSong}{STFangsong.ttf}

\setCJKmainfont{\MySongTi}[
    Path = \fontdir,
    Scale = 1.05,
    AutoFakeBold = false,
    BoldFont = \MySongTiBold,
    AutoFakeSlant = 0.2,
    ItalicFont = \MyKaiTi,
    BoldItalicFont = \MyKaiTi,
]
\setCJKsansfont{\MyHeiTi}[
    Path = \fontdir,
    Scale = 1.05,
    BoldFont = \MyHeiTiBold,
    AutoFakeSlant = 0.2,
]
\setCJKmonofont{\MyHeiTi}[
    Path = \fontdir,
    Scale = 1.05,
    BoldFont = \MyHeiTiBold,
    AutoFakeSlant = 0.2,
]

% ============================================================================
% Scale = 1.05：CJK 正文 10pt → 五号 10.5pt（参考陶哲轩《实分析》实测）。
% Latin/数学字体不受影响，仍为 10pt，与陶书「正文五号 + 数学10pt」一致。
% 各 CJK 族统一缩放以保持行内字号一致；\zihao 标题因乘 1.05 略大于名义值。
\newCJKfontfamily\kaishu{\MyKaiTi}[Path=\fontdir, Scale=1.05, AutoFakeBold=2.5, AutoFakeSlant=0.2]
\newCJKfontfamily\songti{\MySongTi}[Path=\fontdir, Scale=1.05, BoldFont=\MySongTiBold, AutoFakeSlant=0.2]
\newCJKfontfamily\heiti{\MyHeiTi}[Path=\fontdir, Scale=1.05, BoldFont=\MyHeiTiBold, AutoFakeSlant=0.2]
\newCJKfontfamily\fangsong{\MyFangSong}[Path=\fontdir, Scale=1.05, AutoFakeBold=2.5, AutoFakeSlant=0.2]

% 封面标题保留原 SimHei（不随正文黑体替换为思源黑体）
\newCJKfontfamily\coverheiti{SimHei.ttf}[Path=\fontdir, Scale=1.05, AutoFakeBold=2.5, AutoFakeSlant=0.2]

\newcommand\thmheiti{\heiti}

\setCJKfamilyfont{hei2}{\MyHeiTi}[Path=\fontdir, Scale=1.05, BoldFont=\MyHeiTiBold, AutoFakeSlant=0.2]
\newcommand{\heiTwo}{\CJKfamily{hei2}}
\setCJKfamilyfont{sectionfont}{\MyHeiTi}[Path=\fontdir, Scale=1.05, BoldFont=\MyHeiTiBold, AutoFakeSlant=0.2]
\newCJKfontfamily\pffont{\MySongTi}[Path=\fontdir, Scale=1.05, BoldFont=\MySongTiBold, AutoFakeSlant=0.2]

% ============================================================================
% mathrsfs/dsfont/bm were previously loaded here but are unused in the body
% (\mathscr, \mathds, \bm have zero occurrences in chapters/). They conflict
% with unicode-math (now loaded in math-setup.tex), so they are removed.
% \mathfrak (used only in styles/mycommand.sty Lie-algebra macros, themselves
% unused in the body) and \boldsymbol (from amsbsy, 8 uses) are provided by
% base LaTeX / amsmath and survive the removal.

% ============================================================================
% Explicit font-shape fallbacks (silent declared substitutions)
% ============================================================================
% OMS: \textcircled's default accent uses OMS/\rmdefault. Times is not an OMS
% math-symbol font; substitute cmsy so LaTeX does not warn and fall back ad hoc.
\DeclareFontFamily{OMS}{TimesNewRoman.ttf(0)}{}
\DeclareFontShape{OMS}{TimesNewRoman.ttf(0)}{m}{n}{<->ssub * cmsy/m/n}{}
\DeclareFontShape{OMS}{TimesNewRoman.ttf(0)}{b}{n}{<->ssub * cmsy/m/n}{}
\DeclareFontShape{OMS}{TimesNewRoman.ttf(0)}{bx}{n}{<->ssub * cmsy/m/n}{}

% TU sc: neither Times nor Noto Serif SC has OpenType smcp. Map small-caps
% shapes to the matching series upright so \textsc under \bfseries is silent.
\DeclareFontShape{TU}{TimesNewRoman.ttf(0)}{m}{sc}{<->ssub * TimesNewRoman.ttf(0)/m/n}{}
\DeclareFontShape{TU}{TimesNewRoman.ttf(0)}{b}{sc}{<->ssub * TimesNewRoman.ttf(0)/b/n}{}
\DeclareFontShape{TU}{TimesNewRoman.ttf(0)}{bx}{sc}{<->ssub * TimesNewRoman.ttf(0)/b/n}{}
\DeclareFontShape{TU}{NotoSerifSC-Regular.otf(0)}{m}{sc}{<->ssub * NotoSerifSC-Regular.otf(0)/m/n}{}
\DeclareFontShape{TU}{NotoSerifSC-Regular.otf(0)}{b}{sc}{<->ssub * NotoSerifSC-Regular.otf(0)/b/n}{}
\DeclareFontShape{TU}{NotoSerifSC-Regular.otf(0)}{bx}{sc}{<->ssub * NotoSerifSC-Regular.otf(0)/b/n}{}

% ============================================================================
% TEMP-DISABLED for verification (will restore)
% \xeCJKsetcharclass{"2014}{"2014}{1}
% \xeCJKsetcharclass{"2026}{"2026}{1}
\xeCJKDeclareCharClass{CJK}{"2460 -> "24FF}

\let\oldtextcircled\textcircled
\renewcommand{\textcircled}[1]{{\fontspec[Path=\fontdir]{\MyHeiTi}\oldtextcircled{#1}}}
\AtBeginDocument{
    \renewcommand{\textcircled}[1]{{\heiti\oldtextcircled{#1}}}
}
